\section{Related Work}\label{sec:related-work}

Research related to this study can be grouped into four major areas:
deep learning techniques for stock price prediction, the application of
Large Language Models (LLMs) in financial forecasting, retrieval-based
methods for improving the reliability of language models, and hybrid
approaches that combine numerical forecasting with textual reasoning.
Although each of these areas has been widely studied, they have
generally been explored independently. This section reviews the main
contributions of previous research and identifies the limitations that
motivate the proposed work.

\subsection{Deep Learning for Stock Price Prediction}\label{sec:deep-learning}

Long Short-Term Memory (LSTM) networks, introduced by \citet{hochreiter1997}, were designed to overcome the vanishing gradient
problem that affected traditional recurrent neural networks. Their
ability to learn long-term dependencies has made them one of the most
popular architectures for financial time-series forecasting. \citet{nelson2017} demonstrated the effectiveness of LSTMs for predicting stock
price direction using technical indicators such as closing price,
trading volume, and moving averages. Their work also established several
experimental practices, including chronological data splitting, which
continue to be widely adopted in later studies.

Despite these contributions, the evaluation strategy used in many
forecasting studies has an important limitation. Most research compares
model accuracy with a random 50\% baseline, assuming that upward and
downward market movements occur equally often. In reality, stock markets
generally experience more upward movements than downward ones over long
periods. Consequently, a simple strategy that always predicts an upward
movement can achieve accuracy higher than 50\%. In the dataset used in
this study, this naive baseline achieves an accuracy of
\textbf{54.33\%}, which is higher than the performance reported by
several existing forecasting models.

\citet{mehtab2020} emphasized that directional prediction is more
relevant than estimating the exact future stock price because trading
decisions primarily depend on market direction. This perspective is also
adopted in the present work. \citet{gu2020} further showed through a
comprehensive empirical study that stock market predictability is
relatively limited and that small improvements in prediction accuracy
can easily be overstated when inappropriate evaluation methods are used.
Their findings highlight the importance of rigorous statistical testing
before claiming meaningful performance improvements. However, their work
focuses entirely on numerical prediction models and does not consider
the reliability or interpretability of natural language explanations
generated by modern AI systems.

\subsection{Large Language Models in Financial Forecasting}\label{sec:llm-finance}

The growing success of Large Language Models has encouraged researchers
to investigate their potential for financial prediction. \citet{lopezlira2023} were among the first to systematically evaluate whether an
LLM could predict stock market movements using financial news headlines.
Their results indicated that language models contain useful predictive
information and can outperform random guessing in certain scenarios. At
the same time, the authors reported several cases where the model
produced highly confident but incorrect predictions, highlighting the
need for more reliable evaluation methods.

Although their work demonstrated the usefulness of language models, the
comparison was performed only against traditional market benchmarks
rather than against a deterministic numerical forecasting model
evaluated under identical conditions. As a result, it remains unclear
whether the observed predictive capability originates from the language
model itself or from information that could also be captured using
conventional machine learning methods. In addition, the study did not
investigate prediction consistency across repeated executions, even
though LLMs may generate different outputs when presented with the same
input.

\citet{yu2024} proposed a more advanced framework in which the
language model generates a prediction, evaluates its own explanation,
and then refines its response through self-reflection. Their work treats
explanation quality as an important evaluation criterion rather than
simply focusing on prediction accuracy. This idea strongly influenced
the response quality assessment adopted in the present study. However,
the evaluation relies on human annotators and language models acting as
judges. Since language model outputs are inherently non-deterministic
and human evaluations cannot be reproduced exactly, the resulting
quality scores may vary between experiments.

A comprehensive survey by \citet{nie2024}, covering more than fifty
studies on LLM-based financial forecasting, identifies three major
challenges that remain unresolved. These include hallucinated numerical
information when external grounding is unavailable, the inability of
language models to maintain consistent knowledge across multiple
predictions, and weaknesses in numerical reasoning and arithmetic. These
limitations demonstrate the need for more reliable evaluation methods
and motivate the development of grounded and hybrid forecasting systems.

\subsection{Retrieval-Augmented Generation for Reducing Hallucination}\label{sec:rag-hallucination}

Retrieval-Augmented Generation (RAG), introduced by \citet{lewis2020},
was developed to improve the reliability of language models by allowing
them to retrieve relevant information from external knowledge sources
before generating a response. This approach reduces the likelihood of
producing unsupported or fabricated statements and has become widely
adopted for knowledge-intensive question-answering tasks. In these
applications, the correct answer is usually available within the
retrieved documents, making it straightforward to evaluate retrieval
quality.

However, applying RAG to financial forecasting presents additional
challenges. Unlike traditional question-answering tasks, predicting
future stock price movement is not a simple information retrieval
problem. Retrieved documents provide useful background information that
supports reasoning, but they do not directly contain the future market
outcome. Therefore, successful retrieval alone does not guarantee an
accurate prediction.

Another important challenge is numerical reliability. In financial
forecasting, even small numerical mistakes can significantly affect
decision-making. For example, reporting a return of \textbf{-14.57\%} as
\textbf{-0.1457} does not introduce new information, but it changes the
numerical scale and may lead to incorrect investment decisions.
Conventional RAG evaluation mainly focuses on identifying fabricated
facts or unsupported statements and generally does not measure this type
of numerical inconsistency.

Recent studies have extended retrieval techniques to time-series
forecasting. \citet{wang2024} proposed a retrieval-based framework for
improving time-series prediction, while \citet{chen2025} applied
retrieval-augmented language models to large-scale financial
forecasting. Although both studies reported improvements in predictive
performance, neither distinguished between gains in forecasting accuracy
and improvements in numerical correctness. This limitation motivates the
detailed numerical fidelity evaluation performed in the present work.

\subsection{Hybrid Forecasting Systems}\label{sec:hybrid-forecasting}

Hybrid forecasting systems combine multiple predictive approaches to
exploit the strengths of individual models. Several studies reported
improved performance by combining sentiment analysis, deep learning
techniques, and statistical methods.

Most existing approaches rely on fixed weighting strategies that assign
equal importance to all components. Furthermore, they rarely examine the
contribution of individual components separately. The present study
addresses these limitations by introducing a calibrated fusion strategy
and evaluating each component independently.

\subsection{Research Gap}\label{sec:research-gap}

The existing literature demonstrates that each of the major research
directions contributes to stock market forecasting in different ways.
Deep learning models, particularly recurrent neural networks, provide a
strong baseline for analysing financial time-series data. Large Language
Models have shown the ability to extract useful information from
financial text, while Retrieval-Augmented Generation improves the
reliability of language model responses by grounding them in external
evidence. In addition, several studies have reported that hybrid systems
combining multiple models often achieve better overall performance than
individual approaches.

Although these findings are promising, an important research question
remains unanswered. Previous studies typically evaluate only one or two
of these components at a time, making it difficult to determine the
individual contribution of language reasoning, retrieval grounding, and
model fusion. Since these techniques are introduced simultaneously, it
is unclear which component is responsible for the observed improvement
and whether the improvement is consistent under identical experimental
conditions.

To address this limitation, the present study performs a controlled
comparison in which four forecasting systems are evaluated using the
same dataset, identical input samples, and a common evaluation
framework. This experimental design allows the effect of each
engineering layer to be measured independently while minimizing the
influence of other variables.

Table 1 compares six representative studies with the capabilities
required for this research. The comparison highlights that, although
previous work has advanced individual aspects of financial forecasting,
no existing study evaluates all four system configurations under
identical conditions while simultaneously measuring forecasting
accuracy, numerical reliability, prediction stability, hallucination,
and response quality. This methodological gap forms the primary
motivation for the proposed research.

\begin{table}[H]
\centering
\scriptsize
\setlength{\tabcolsep}{2.8pt}
\resizebox{\textwidth}{!}{%
\begin{tabular}{@{}>{\raggedright\arraybackslash}p{4.7cm}ccccccc@{}}
\toprule
\textbf{Capability required by this study} & \textbf{Nelson (2017)} & \textbf{Gu et al. (2020)} & \textbf{Lopez-Lira \& Tang (2023)} & \textbf{Lewis et al. (2020)} & \textbf{Yu et al. (2024)} & \textbf{Fatouros (2024)} & \textbf{This work} \\
\midrule
C1 Deterministic ML arm scored on identical rows & Yes & Yes & No & No & No & Yes & Yes \\
C2 Same LLM run with and without retrieval & n/a & n/a & No & Partly & No & No & Yes \\
C3 Fusion calibrated to each arm's own decision point & n/a & No & n/a & n/a & No & No & Yes \\
C4 Regime-split baseline, not a 50\% coin flip & No & Partly & No & n/a & No & No & Yes \\
C5 Mechanical hallucination audit of numeric claims & n/a & n/a & No & Partly & No & n/a & Yes \\
C6 Run-to-run determinism measured, not assumed & n/a & n/a & No & No & No & n/a & Yes \\
C7 Answer quality scored without a human or LLM judge & No & No & No & No & No & No & Yes \\
C8 Attainable-ceiling analysis for the task & No & Partly & No & No & No & No & Yes \\
\bottomrule
\end{tabular}}
\caption{Capability comparison against six studies. ``n/a'' marks a capability that is not applicable to that study's design rather than a shortcoming of it.}
\label{tab:capability-comparison}
\end{table}

The final column of \textbf{Table 1} shows that none of the reviewed
studies is fundamentally incomplete or incorrect. Instead, each study
was designed to answer a specific research question using an appropriate
methodology. For example, \citet{nelson2017} and \citet{gu2020}
focused on evaluating numerical forecasting models through rigorous
experimental procedures but did not examine the reliability of generated
explanations. \citet{lopezlira2023} investigated the forecasting
capability of a Large Language Model without comparing it against an
equivalent deterministic machine learning model. \citet{lewis2020}
demonstrated the effectiveness of Retrieval-Augmented Generation in
knowledge-intensive tasks where the correct answer is available within
the retrieved documents. \citet{yu2024} emphasized explanation
quality, but their evaluation depended on human reviewers and
language-model-based judges, making the assessment difficult to
reproduce consistently. Similarly, \citet{fatouros2024} proposed a
hybrid forecasting framework by combining multiple prediction signals;
however, the fusion relied on fixed weights and did not include a
generative language model as an independent component.

Based on this review, the key research gap becomes clear. To the best of
our knowledge, no previous study has evaluated a standalone LLM, a
retrieval-enhanced LLM, a deterministic sequence model, and a calibrated
hybrid of these models using the same dataset and identical evaluation
conditions. Furthermore, existing studies rarely assess these systems
simultaneously in terms of prediction accuracy, numerical reliability,
response consistency across repeated executions, hallucination
behaviour, and overall response quality.

The present study addresses this methodological gap by introducing a
controlled evaluation framework rather than proposing a completely new
forecasting architecture. Well-established components, including a
two-layer LSTM, TF-IDF-based retrieval, and a calibrated log-odds fusion
strategy, are intentionally selected to ensure that the comparison
focuses on the contribution of each engineering layer. Consequently, the
main novelty of this work lies in the evaluation methodology, which
enables the individual effects of language reasoning, retrieval
grounding, and model fusion to be measured under identical experimental
conditions.
